\section{The Thesis: In Natural Units $\alpha = 1$}

\textbf{[K]} The strongest statement about the fine structure constant is the
simplest: in natural units
\[
  \alpha = 1 .
\]
The famous value $1/137{.}036$ is \emph{not} a mystery of nature, but the
bookkeeping that arises when converting to SI units -- just as
$c = 3\cdot10^8$ is the bookkeeping of the metre-second choice.

\textbf{This is not an invention of this theory.} Setting $\alpha$ to one is a
\emph{known} unit convention: one normalises the charge so that
$e^2 = 4\pi$ (rationalised Heaviside-Lorentz units with
$\varepsilon_0 = \mu_0 = 1$). In these units $\alpha = 1$ by construction.
The theory adopts this convention; it does not re-justify it.

\textbf{But the reason is different from $c, \hbar, k_B$.} For
$c, \hbar, k_B$, the conversion factors become one because time-mass duality
identifies their dimensions (A085). $\alpha$ is dimensionless; this justification
does \emph{not} apply here. The convention $\alpha = 1$ rests on the
electromagnetic charge normalisation -- a separate reason having nothing to do
with time-mass duality and therefore handled here separately, not under the
natural-units topic.

\section{How $\alpha = 1$ Is Derived}

\textbf{[B]} The derivation runs through the charge normalisation and
Maxwell duality, not through an axiom. In SI:
\[
  \alpha = \frac{e^2}{4\pi\varepsilon_0\hbar c}.
\]
With $\hbar = c = 1$ and the natural choice $\varepsilon_0 = 1$, the charge
normalisation $e^2 = 4\pi$ follows from $\alpha = 1$. The Maxwell relation
$c^2 = 1/(\varepsilon_0\mu_0)$ with $c = 1$ and $\varepsilon_0 = 1$ then gives
$\mu_0 = 1$. Substituting confirms:
\[
  \alpha = \frac{e^2}{4\pi\varepsilon_0\hbar c} = \frac{4\pi}{4\pi\cdot1\cdot1\cdot1} = 1 .
\]

\textbf{[K]} The core is electromagnetic duality
$\tfrac{1}{\varepsilon_0 c} = \mu_0 c$, which is perfectly satisfied in natural
units ($1 = 1$). $\alpha = 1$ is the expression of this balance between
electric ($\varepsilon_0$) and magnetic ($\mu_0$) coupling.

\textbf{Honest assessment.} The step itself is standard: the charge
normalisation $e^2 = 4\pi$ is the rationalised convention in which
$\alpha = 1$ comes out. What is non-standard is the \emph{interpretation}:
that the SI value $1/137$ is thereby not a fundamental mystery of the coupling
but measures the historical charge unit. Not everyone shares this interpretation
-- the classical Duff position holds $\alpha$ to be absolute on account of its
dimensionlessness. The theory's actual contribution lies not here, however, but in
the next section: that the SI value not only measures charge units but, via
$\alpha = \xi E_0^2$, carries the geometric parameter $\xi$.

\textbf{This removes the mystification.} The question ``why exactly $1/137$?''
is a question for the SI system, not for nature -- of the same kind as ``why
exactly $299{,}792{,}458$ metres per second?''. The numerical value measures the
historical charge unit, not a property of the coupling itself.

\section{What Remains: $\xi$ and One Scale}

\textbf{[K]} When all constants are one in natural units, the only genuine
content remaining is the single dimensionless parameter $\xi$ and one scale.
The SI value of $\alpha$ carries precisely these two quantities:
\[
  \alpha_\text{SI} = \xi\cdot\Bigl(\frac{E_0}{1\,\text{MeV}}\Bigr)^2 .
\]

Here $\alpha$ is no longer the native one but its SI dress. In this dress one
sees that $\alpha$ and $\xi$ encode the same quantity, connected by the scale
$E_0$: $\xi$ is the geometric reading (packing ratio, A020), $\alpha_\text{SI}$
the electromagnetic one. Two sides of one coin, one degree of freedom -- just as
mass, energy, temperature, and inverse time are one quantity in four dresses (A085).

\section{Why the Measured Value Is Irrelevant}

\textbf{[K]} The theory works with ratios (A080). Ratios are unit-free and depend
only on the geometry. The absolute value of $\alpha$ -- $1/137{.}036$ -- is an
absolute value and therefore belongs to the bookkeeping layer, not the core.

\textbf{[K]} The \emph{structural relations} at issue are already known and
confirmed in SI units and in the Standard Model. The theory does not invent them,
it orders them: $\alpha = \xi E_0^2$ connects known quantities. What it asserts
is the \emph{structure} of this connection -- that $\alpha$ and $\xi$ are the same
degree of freedom -- not a new numerical value for $\alpha$.

\textbf{[B]} It follows how the numbers in the following sections are to be read:
\emph{not} as a test the theory passes or fails. They show the precision with
which the bookkeeping reproduces the SI value -- a statement about the bridge
constant, not about the structure. As long as one works with ratios, they are
irrelevant.

\textbf{[STIPULATION]} The test of the theory lies elsewhere. $\alpha$ is not one
of the four prediction quantities of the series (A010); those are ratios and
structural statements ($m_\tau$, the $\xi$-order of the CHSH deviation, $D_f$,
$L_0$). Whether $\alpha$ approaches $1/137$ to three or five digits does not test
the theory -- it tests the SI conversion.

\section{Why the Circularity Is Not a Weakness}

\textbf{[STIPULATION]} As soon as one works in SI dress, the system is circular,
and rightly so. Every system that fixes its own constants is circular -- the only
question is where one begins. (In natural units the question does not arise at all:
there $\alpha = 1$ and nothing to determine.)

One can start from $\xi$ (geometric perspective): then $\xi$ is primary and
$\alpha = \xi E_0^2$ derived. One can start from $\alpha$ (electromagnetic
perspective): then $\alpha$ is primary and $\xi = \alpha/E_0^2$ derived. Both
routes give the same values for all other constants. They are mathematically
equivalent.

\textbf{[B]} The question ``can $\alpha$ be derived from $\xi$ without circularity?''
is therefore wrongly posed. In a closed system there is no exterior from which
one could derive. There is an entry point one declares, and the rest that follows
from it. Classical physics conceals this by placing several constants side by side
as ``fundamental''; this theory makes explicit that only one entry point is needed.

\textbf{A mistake would still be possible here,} namely presenting the entry point
as a derivation from nothing. Whoever substitutes $E_0 = \sqrt{\alpha/\xi}$ and
then announces ``exactly the measured value of $\alpha$'' confuses reading consistency
with proof. The $E_0$ determination avoids exactly this mistake, because it does
\emph{not} run via $\sqrt{\alpha/\xi}$: the exact $E_0$ is reached by Route~2
independently of $\alpha$ (previous section), and the agreement of the two
$\alpha$-free routes makes the $\alpha$ point over-determined rather than merely
consistent. The circularity that always threatens in a closed system is therefore
broken at this point by two separate approaches.

\section{The Scale \texorpdfstring{$E_0$}{E0} as Entry Point}

\textbf{[K]} What separates $\xi$ from $\alpha$ is only the scale $E_0$. For it
there are \emph{two mutually independent} routes, both not using $\alpha$
and confirming each other.

\textbf{[K]} \emph{Route~1 (empirical).} The geometric mean of the two
lightest lepton masses,
\[
  E_0 = \sqrt{m_e\,m_\mu} = 7{.}348\ \text{MeV},
\]
gives when substituted $\alpha = \xi\cdot(7{.}348)^2 = 1/138{.}9$, thus $-1{.}35\,\%$
from the measured value.

\textbf{[B]} \emph{Route~2 (geometric).} A construction from $\xi$ and
$m_\mu$ alone,
\[
  E_0^2 = 4\sqrt2\;\frac{m_\mu}{\xi^{p}},
\]
gives $E_0 = 7{.}398$ MeV -- \emph{without} using the geometric mean,
\emph{without} $m_e$, and \emph{without} back-calculation from $\alpha$. Substituted,
this hits $\alpha$ to $5\cdot10^{-6}$.

\textbf{[K]} Decisive is that the two routes \emph{meet} without knowing each other.
Their ratio is
\[
  \frac{7{.}398}{7{.}348} = 1{.}00682 = K_\text{frak}^{-1/2}
  \quad\text{to } 9\cdot10^{-5},
\]
with the same fractal factor $K_\text{frak} = 1-100\xi$ derived in A040
and acting in the corpus at an entirely different point. Since $E_0$ enters
$\alpha$ quadratically, the $K_\text{frak}^{1/2}$ bridge at $E_0$-level corresponds
to a full $K_\text{frak}$ correction at $\alpha$-level. Two separate derivations
thus meet to $8\cdot10^{-5}$, and the depth of the correction is not free but
structurally fixed.

\textbf{[K]} The $\alpha$ point is thereby \emph{not a calibration but
over-determined}. The exact $E_0 = 7{.}398$ is not ``formed backwards from
$\alpha$'' but reached forward by Route~2; and that precisely the derived
$K_\text{frak}$ closes the gap between the two routes is a cross-sector consistency,
not a fit. Because Route~2 does not use $m_e$, the agreement with Route~1 is even
a genuine statement about $m_e$: the remaining residual of $8\cdot10^{-5}$ is too
large for mass measurement noise and is thus a test of the pole mass against
the geometry.

\section{The Two Paths Also Confirm \texorpdfstring{$K_\text{frak}$}{K\_frak}}

\textbf{[K]} The over-determination says more than has been exploited so far.
Neither Route~1 nor Route~2 uses $K_\text{frak}$ anywhere: Route~1 is the
geometric mean of two measured masses, Route~2 a construction from $\xi$
and $m_\mu$. Their ratio is therefore a \emph{measurement} of $K_\text{frak}$
that does not presuppose the factor:
\[
  K_\text{frak} \;=\; \Bigl(\frac{E_0^{\,(2)}}{E_0^{\,(1)}}\Bigr)^{-2}
  \;=\; \Bigl(\frac{7{.}398}{7{.}348}\Bigr)^{-2} \;=\; 0{.}9862 .
\]

\textbf{[K]} This value read from the two routes hits the factor independently
derived in A040:
\[
  K_\text{frak} = 1 - 100\xi = \frac{74}{75} = 0{.}98667
\]
to $9\cdot10^{-5}$. This is remarkable because the two origins have nothing to do
with each other: A040 obtains $K_\text{frak}$ from the cumulative RG run and the
consistency condition on the effective fractal dimension; the $\alpha$ sector
obtains \emph{the same} value from the difference between two energy-scale routes.
A factor originating from an RG argument that simultaneously closes the gap of two
independent $E_0$ determinations is not freely chosen; it is fixed from two sides.

\textbf{[B]} $K_\text{frak}$ thus has a \emph{second witness} here. The $\alpha$
sector confirms not only the value of $E_0$ and $\alpha$, but via the route
difference also that $K_\text{frak} = 74/75$ has the correct value -- without
inserting this factor anywhere. What in A040 is a derived constant is here
additionally \emph{measured}.

\section{A Coincidence That Must Be Resisted}

\textbf{[K]} It is tempting to search for a \emph{more elegant} expression for $E_0$
than the mass route, and there is a striking one: Euler's number squared,
\[
  \mathrm{e}^2 = 7{.}389\ \text{MeV}\ (?),\qquad
  \alpha = \xi\cdot(7{.}389)^2 = \frac{1}{137{.}37}\quad(-0{.}24\,\%),
\]
which hits the SI value even better than the mass route. This expression should be
\emph{resisted}, for a principled reason: it equates a pure number (Euler's number)
with an energy in the arbitrary unit MeV. One MeV is a human convention; no pure
number can fix a physical scale \emph{in} MeV without the agreement depending on
the unit choice. In a different energy unit a different number would appear, and
the coincidence would vanish.

\textbf{[K]} This is exactly the pattern the corpus warns against elsewhere
(A105): a number fitted backwards to an elegant expression without forward content.
That $\mathrm{e}^2$ hits closer than $\sqrt{m_e m_\mu}$ does \emph{not} make it
the better candidate -- the mass route carries a physical reason (the lepton masses),
the Euler coincidence carries none. It is recorded here only to explicitly
\emph{not} book it as an anchor. The elegant expression is not needed anyway:
the exact $E_0$ is already hit by Route~2 with forward content, while $\mathrm{e}^2$
merely holds a number to a unit.

\section{Verification Script}

\begin{itemize}
  \item The equivalence of both perspectives ($\alpha = \xi E_0^2$ and
    $\xi = \alpha/E_0^2$ deliver identical values); the $\alpha$-free
    harmonic route with $\kappa = 4$; the two numerical routes; and the explicit
    check that substituting $\sqrt{\alpha/\xi}$ is a consistency and not an
    independent proof:
    \path{python/A_Serie_Skripte/a130_alpha_kette.py}
\end{itemize}

\section{Open Questions}

\textbf{First, the exact $E_0$ is $\alpha$-free and over-determined -- not open.}
The exact $E_0 = 7{.}398$ is reached by Route~2 from $\xi$ and $m_\mu$, without
$\alpha$ and without $m_e$; it agrees with the mass route (Route~1) up to the
derived factor $K_\text{frak}^{-1/2}$. The earlier classification that the exact
value was formed only backwards from $\alpha$ is thereby superseded. The Euler
coincidence $E_0 = \mathrm{e}^2$ remains explicitly \emph{not} a derivation and
is not needed.

\textbf{Second, the absolute anchor remains tied to a measured mass.} Route~2 fixes
$E_0$ from $\xi$ and $m_\mu$; the absolute scale therefore enters via \emph{one}
measured mass anchor. A determination of the exact $E_0$ from purely dimensionless
geometry without any mass anchor does not exist -- but this is the general structure
of the SI bridge (A080), not a deficiency specific to the $\alpha$ sector. Since
the theory tests via ratios, the remaining residual carries no weight for its
evaluation.

\textbf{Third, a lower bound applies.} The relative description
$\alpha \leftrightarrow \xi$ is secured down to the sub-Planck bound $L_0 = \xi\,\ell_P$
(A210); below this the theory makes no statement.

\section{Supersedes}

This document subsumes: Doc.~011 (fine structure constant), Doc.~043 (resolution
of constants), Doc.~044, Doc.~101 (circularity of constants, core thesis: one
degree of freedom, $\xi$ and $\alpha$ equivalent). Incorporated: P40.

\section{Use of AI Tools}

This document was created with the assistance of AI tools. The questions,
stipulations, and all substantive decisions come from the author; the tools
formulate, compute, and create verification scripts.
Responsibility for content and errors rests entirely with the author.
Full wording of the rules in A010.
